%% Chapter 5 - Implementation, Integration and Testing
%% Planned sections (master structure): 5.1 Development environment and
%% technology stack; 5.2 Component implementation; 5.3 Integration with the
%% other components; 5.4 Verification and software testing; 5.5 Deployment and
%% operation.
\chapter{Implementation, Integration and Testing}\label{ch:implementation-testing}

\section{Development Environment and Technology Stack}\label{sec:dev-environment}

The IRE is implemented as a modular Python service with a FastAPI interface
and a clean internal separation between the core stages (S1--S7), the guard
middleware (M1--M5) and the language-model adapter (Chapter~\ref{ch:product-design}).
This section records the verified technology stack and the development
environment in which the system was built and measured.

\subsection*{Runtime and deployment platform}

The service is written in Python~3 and runs as a local daemon under Uvicorn on
port 8001. The development launcher (\texttt{scripts/run\_dev.py}) targets
Python 3.11+ and performs pre-flight checks for the \texttt{.env} file, Ollama
connectivity on port 11434, the configured model's availability and the
development server. The evaluation and QA environments for this dissertation
were exercised under Python 3.9.13 (the retained virtual environment), which
is below the recommended version but within the declared compatibility window.
Table~\ref{tab:stack} lists the verified stack, pinned where a version is
recorded in \texttt{requirements.txt}.

\begin{table}[h]
\centering
\small
\caption{Verified technology stack.}\label{tab:stack}
\begin{tabular}{@{}lll@{}}
\toprule
\textbf{Layer} & \textbf{Technology} & \textbf{Purpose}\\
\midrule
Language & Python 3.11 (3.9.13 used) & Implementation platform\\
API framework & FastAPI 0.115.0, Uvicorn 0.30.6 & HTTP surface, ASGI server\\
Configuration & pydantic-settings 2.4.0, python-dotenv & \texttt{.env}-driven settings\\
Schemas & Pydantic v2 ($\ge 2.9$, $<3$) & Typed model validation\\
Inference client & ollama 0.3.3, httpx 0.27.2 & Local Ollama decoding\\
Rate limiting & slowapi 0.1.9 & 30 requests/min per caller\\
Logging & structlog 24.4.0 & Structured JSON audit\\
Embeddings & sentence-transformers 3.0.1 & Cache similarity, consistency\\
Testing & pytest 8.3.2, pytest-asyncio 0.23.8 & Unit + async integration\\
\bottomrule
\end{tabular}
\end{table}

Two declared dependencies are flagged for accuracy: \texttt{spacy} 3.7.6 is
listed in \texttt{requirements.txt} and probed by the launcher, and
\texttt{jsonschema} 4.23.0 is listed in the SRS, but the implemented validation
paths in \texttt{src/} use the Pydantic runtime, not these packages. They are
treated in this dissertation as declared-but-unused rather than as part of the
runtime (Chapter~\ref{ch:product-design}).

\subsection*{Source layout and configuration}

\begin{sloppypar}
The service is organised under
\texttt{src/}: \texttt{api}, \texttt{audit}, \texttt{inference},
\texttt{integration}, \texttt{middleware}, \texttt{pipeline},
\texttt{preprocessing}, \texttt{schemas}, \texttt{utils} and
\texttt{validation}. Central configuration lives in \texttt{config/}\hspace{0pt}\texttt{settings.py}
with the feature, RBAC, and enforcement policies in \texttt{features.yaml},
\texttt{feature\_}\hspace{0pt}\texttt{flags.py}, \texttt{rbac\_}\hspace{0pt}\texttt{policy.py} and
\texttt{enforcement\_}\hspace{0pt}\texttt{policy.py}. Runtime constants are not decided by the
code: they are loaded from the environment and the YAML feature flags at
startup, so every experimental constant is reproducible from the committed
configuration (Section~\ref{sec:reproducibility}).
\end{sloppypar}


\section{Component Implementation}\label{sec:component-implementation}

This section describes how the design of Chapter~\ref{ch:product-design} is
realised in code. The descriptions are based on the implementing modules under
\texttt{src/} and \texttt{config/} as committed.

\subsection*{Pipeline orchestration}

\texttt{src/pipeline/ire\_pipeline.py} instantiates every stage and middleware
and defines their execution order. Each middleware and stage is wrapped in its
own exception boundary so that a failure in non-critical middleware degrades
gracefully instead of aborting the request (Section~\ref{sec:requirements}).
The order is: M1 adversarial scan, M2 pre-inference RBAC, M3 session-context
injection, M4 semantic-cache lookup, then the core S1--S7 with S6B, then M5
post-inference RBAC, cache store, v2-response assembly, sub-intent enrichment,
session update and session-info enrichment. Every stage and middleware that
can reject a request does so with a classified error carrying the stage name,
the divergence field and its accumulated \texttt{ValidationFinding} entries.

\subsection*{Core stages}

\begin{itemize}
\item \textbf{S1 InputNormalizer} (\texttt{preprocessing/input\allowbreak\_normalizer.py})
applies NFKC Unicode normalisation, strips control characters and thinking
tags, collapses whitespace and standardises a jargon table (for example
\texttt{syn scan} $\rightarrow$ \texttt{SYN stealth scan},
\texttt{pwn} $\rightarrow$ \texttt{exploit}).
\item \textbf{S2 AliasResolver} (\texttt{preprocessing/alias\allowbreak\_resolver.py})
resolves a 22-entry \texttt{data/alias\_map.json} of exploit surnames, slang
and tool terms to canonical identifiers before inference.
\item \textbf{S3 OllamaInferenceEngine} calls the local model with a few-shot,
strict-JSON directive, a per-stage timeout (60 s), a 2048-token cap and a
small response cache
(\texttt{inference/ollama\_inference\_engine.py}).
\item \textbf{S4 JSONParser} (\texttt{validation/json\_parser.py}) recovers the
JSON payload, stripping code fences, markdown artefacts and hard bounds on
extraction.
\item \textbf{S5 SchemaValidator} (\texttt{validation/schema\allowbreak\_validator.py})
enforces the typed \texttt{IntentSchema} with Pydantic and maps failures to
hallucination classes (fabricated CVE, target-type mismatch, ungrounded
confidence, fabricated parameter).
\item \textbf{S6 RegexValidator} re-verifies IP/CIDR/CVE/port/URL entities and
rejects shell metacharacters (\texttt{validation/regex\_validator.py}).
\item \textbf{S6B NetworkArchitectureValidator} checks engagement-scope
containment, CIDR sanity (widest prefix 30, narrowest-32 scan targets) and
target-type consistency (\texttt{validation/network\_validator.py}).
\item \textbf{S7 ScopeGuard} (\texttt{validation/scope\_guard.py}) enforces
RFC~1918 membership~\cite{rfc1918}, the engagement mode (strict/warn/disabled)
and re-scans the raw input as a final injection gate.
\end{itemize}

\subsection*{Guard middleware}

\begin{itemize}
\item \textbf{M1 AdversarialDetector}
(\texttt{middleware/adversarial\_detector.py}) implements the three-layer
detector of the design: a lexical layer scoring known injection phrases and
control tokens, a structural layer scoring encoded payloads and script-like
constructs, and a scoring layer that auto-blocks at a cumulative score of at
least 0.95 and flags at 0.7.
\item \textbf{M2/M5 RbacGuard} is driven by \texttt{config/rbac\_policy.py},
which declares the role hierarchy (viewer, analyst, operator, admin), the
intent-permission matrix, the three sensitive intents and the restricted-role
set (\texttt{viewer}); unknown roles coerce to \texttt{viewer}.
\item \textbf{M3 SessionContextStore}
(\texttt{middleware/session\_context.py}) tracks per-session turn history with
a TTL, injects a memory window and carries the last target forward.
\item \textbf{M4 SemanticCache} (\texttt{middleware/semantic\_cache.py}) serves
exact-hash and embedding-similarity hits and is invalidated with the CVE
grounding.
\item \textbf{sub-intent classifier}
(\texttt{middleware/sub\_intent\_classifier.py}) enriches a valid response with
a hierarchical sub-intent that drives tool selection.
\end{itemize}

\subsection*{Contract generation, audit and metrics}

\begin{sloppypar}
The downstream \texttt{PlannerContract} is synthesised by
\texttt{src/pipeline/contract\_builder.py}: sub-intent selects the tool (for
example \texttt{NETWORK\_SCAN.SYN\_STEALTH} $\rightarrow$ \texttt{nmap -sS}),
modifiers map to flags, priorities are assigned, and exploitation flows are
marked for human approval. The audit trail in
\texttt{audit/}\hspace{0pt}\texttt{audit\_logger.py} persists every finding to
an append-only SQLite database (\texttt{data/audit\_log.db}), hashing operator
content (SHA-256) while keeping intent, role, latency and outcome as
structured fields. \texttt{utils/}\hspace{0pt}\texttt{metrics\_collector.py}
records latency percentiles, per-stage error counts and intent distribution
for the \texttt{/metrics} endpoint.
\end{sloppypar}


\section{Integration with the Other Components}\label{sec:integration}

The IRE is Component~1 of the NeuroShell framework and integrates with
Components 2--4 under feature flags, so the framework can be turned on without
changing the parsing behaviour. \texttt{src/}\hspace{0pt}\texttt{integration/}\hspace{0pt}\texttt{planner\_client.py}
forwards validated contracts to the Dynamic Planner (Component~2, port 8002)
when \texttt{planner\_enabled} is true;
\texttt{src/}\hspace{0pt}\texttt{integration/}\hspace{0pt}\texttt{executor\_client.py}
targets the AEERE executor (Component~3, port 8003) and the AVAE vulnerability
analysis (Component~4, port 8004) under their respective flags. Clients
authenticate with a shared API key over HTTP and a bounded timeout.

Two integration modes are implemented in the API surface. \texttt{POST /parse}
returns the validated intent and pushes a background task to the planner when a
contract is actionable; \texttt{POST /execute} runs the synchronous C1-to-C2
chain and returns both outputs, persisted in
\texttt{utils/}\hspace{0pt}\texttt{chain\_cache.py} so
an identical input short-circuits both services. \texttt{POST /parse/plan}
returns the planner-ready contract directly. Downstream consumers receive
\texttt{PlannerContract} objects only, never free-form model output; the
Component~2 safety gates (RFC~1918, dangerous-pattern checks) remain a second,
independent enforcement point (Section~\ref{sec:component-architecture}).

The full API surface, implemented in \texttt{src/api/main.py}, is:
\begin{itemize}
\item GET \texttt{/health};
\item GET \texttt{/metrics};
\item POST \texttt{/parse};
\item POST \texttt{/parse/plan};
\item POST \texttt{/execute};
\item GET \texttt{/session/}\texttt{\{id\}} and DELETE \texttt{/session/}\texttt{\{id\}};
\item GET \texttt{/admin/rbac/}\texttt{\{role\}};
\item GET \texttt{/admin/features} and POST \texttt{/admin/features/reload};
\item GET \texttt{/admin/models} and POST \texttt{/admin/models/switch};
\item GET \texttt{/admin/enforcement-policy};
\item GET \texttt{/admin/escalations} and POST
\texttt{/admin/escalations/}\hspace{0pt}\texttt{\{id\}/}\hspace{0pt}\texttt{resolve};
\item GET \texttt{/admin/audit/summary};
\item GET \texttt{/admin/chain-cache} and DELETE \texttt{/admin/chain-cache};
\item GET \texttt{/engagement/scope}.
\end{itemize}
Requests are authenticated with an API key and rate-limited to 30 per minute
per caller; command length is capped in the API.


\section{Verification and Software Testing}\label{sec:testing}

\subsection*{Strategy}

The verification strategy distinguishes unit-level determinism from
integration-level wiring. Unit tests exercise every guard path, every cell of
the RBAC matrix, both semantic-cache hit and miss paths, all six hallucination
classes, and the metrics calculations, using a live-free harness in which the
LLM is mocked so the suite is deterministic (no live model dependency in the
unit or integration paths). Integration tests exercise the HTTP surface with
an \texttt{httpx.AsyncClient}: the asynchronous parse flow, the adversarial-to-
alias sequence, the full-stack guarded pipeline, and the audit and escalation
queue lifecycle. The suite is the gate that must pass before any measurement
in the results chapter is accepted (Section~\ref{sec:reproducibility}).

\subsection*{Re-derived results}

The figures in this section were re-derived by running the committed suite
directly for this dissertation (session 2026-10-07), because the retained
coverage artefact was internally inconsistent with the current source
(Chapter~\ref{ch:methodology}).

\begin{itemize}
\item \textbf{Suite.} 506 collected tests, 506 passed, 0 failed in the full
run (111.76\,s; 49.7--55.3\,s across runs). The split is 434 unit tests and 72
integration/async tests, an increase from the 482 tests (434 unit + 48
integration) documented on 2026-08-26 because additional integration cases were
added to the API, full-stack and audit suites.
\item \textbf{Coverage.} 76\% statement coverage across \texttt{src/} (2421
statements, 582 missed), with 17 modules at 100\%. The safety-critical core is
well covered: adversarial detection 99\%, schema validation 99\%, session
context 96\%, audit logger 94\%, semantic cache 94\%, sub-intent classifier
89\%. The weakest areas are the integration clients (planner 30\%, executor
16\%), \texttt{contract\_builder.py} (30\%) and \texttt{chain\_cache.py}
(24\%), which are orchestration paths rather than safety invariants.
\item \textbf{Static analysis.} \texttt{ruff} reports 593 findings, 100\%
style and modernisation, with no correctness-level (E9) errors and no fatal
findings.
\end{itemize}

\subsection*{Evaluation harness}

The evaluation harness applies the procedure of
Section~\ref{sec:procedure} to the golden data set, computes the metrics of
Section~\ref{sec:measures}, and regenerates every table used in
Chapter~\ref{ch:results} from a single, manifest-pinned command
(Section~\ref{sec:reproducibility}). The harness records per-record decisions
so the reported figures can be reproduced without re-invoking the model, and
it asserts the gate condition (suite green) before accepting a measurement.


\section{Deployment and Operation}\label{sec:deployment-operation}

\subsection*{Development deployment}

A single command (\texttt{python scripts/run\_dev.py}) performs the pre-flight
checks and starts Uvicorn on port 8001 with hot-reload for development. The
service binds to 0.0.0.0 for the downstream components on the device; no
external service is contacted. Health and readiness are exposed at
\texttt{/health}, which reports Ollama connectivity, the active model, the
seven-stage pipeline and the middleware enablement map.

\subsection*{Containerised deployment}

A \texttt{Dockerfile} provides a reproducible deployment on Python 3.11-slim,
installing the pinned requirements and packaging \texttt{src/}, \texttt{config/}
and the alias map; a \texttt{docker-compose.yml} (documented in the SRS)
orchestrates the service alongside the Ollama runtime, and the image exposes a
health check against \texttt{/health}. The container runs a single worker by
design, because the in-memory cache and session store are per-process state
(Section~\ref{sec:component-architecture}).

\subsection*{Operational administration}

The admin surface supports runtime reconfiguration without redeploys: model
listing and switching, feature-flag inspection and hot reload, live
enforcement-policy summaries, escalation resolution, audit-trail summaries and
engagement-scope inspection. Rate limiting and command-length caps protect the
model backend, and session/cache storage can be inspected and flushed through
the API. Deployment properties are traceable to the trust boundary described
in Chapter~\ref{ch:product-design} and the threat-model treatment of
Chapter~\ref{ch:methodology}.